\subsection{Quark and lepton flavors with common modulus $\tau$ in $A_4$ modular symmetry --- arXiv:2005.00775}

\textbf{Authors:} Hiroshi Okada, Morimitsu Tanimoto

\paragraph{主な主張}
共通のモジュライ $\tau$ をもつモジュラー $A_4$ 模型で、重み2・4・6の形式を用いると固定点 $\tau=i$ の近傍でクォークとレプトンの混合を同時に記述できる領域があることを示した\cite{Okada:2020rjb}。

\paragraph{新規性}
両セクターで許される $\tau$ 領域を照合し、固定点からの変位によるCP破れと、ニュートリノ質量和による模型の検証可能性を結びつけた。

\paragraph{位置付け}
単一モジュライによるクォーク・レプトン共通記述の可能性と制約を調べる具体例である。
